\documentclass[12pt]{article}
\usepackage{geometry}
\geometry{margin=1in}

% ----------------------------------------------------------
%  Linux Penguin Theme — Based on Recommended Colors
% ----------------------------------------------------------

% --- COLOR SETUP ---
\usepackage{xcolor}

\definecolor{cream}{RGB}{246,242,219}
\definecolor{teamred}{RGB}{209,57,64}
\definecolor{navy}{RGB}{40,46,87}
\definecolor{softwhite}{RGB}{250,250,230}
\definecolor{accentyellow}{RGB}{240,200,60}

% Extra utility colors
\definecolor{muted}{RGB}{120,125,140}
\definecolor{darkcream}{RGB}{230,225,205}

% ----------------------------------------------------------
%  FONT SETUP (Optional but Recommended)
%  Works best with XeLaTeX or LuaLaTeX
% ----------------------------------------------------------
\usepackage{fontspec}

% Main document font (Google fonts — upload to Overleaf)
% If you don't upload fonts, comment these two lines.
\setmainfont{SourceSans3}[
  Path=./,
  ItalicFont  = *-Italic,
  BoldFont    = *-Bold
]

\newfontfamily\headingfont{ScienceGothic}[
  Path=./,
]

% ----------------------------------------------------------
%  SECTION STYLE
% ----------------------------------------------------------
\usepackage{titlesec}

% Section
\titleformat{\section}
  {\Large\headingfont\bfseries\color{navy}}
  {\thesection}
  {1em}
  {\rule{\linewidth}{1pt}\vspace{0.5em}\\}

% Subsection
\titleformat{\subsection}
  {\large\headingfont\bfseries\color{teamred}}
  {\thesubsection}
  {0.8em}
  {}

% Subsubsection
\titleformat{\subsubsection}
  {\bfseries\color{navy}}
  {\thesubsubsection}
  {0.6em}
  {}

% ----------------------------------------------------------
%  PAGE STYLE
% ----------------------------------------------------------
\usepackage{fancyhdr}
\pagestyle{fancy}

\fancyhf{} % clear all

\lhead{\textcolor{navy}{\headingfont \courseTitle}}
\rhead{\textcolor{teamred}{\thepage}}
\renewcommand{\headrule}{\color{teamred}\hrule height 1pt}

% ----------------------------------------------------------
%  HYPERLINK STYLE
% ----------------------------------------------------------
\usepackage{hyperref}

\hypersetup{
  colorlinks = true,
  linkcolor  = teamred,
  urlcolor   = navy,
  citecolor  = accentyellow
}

% ----------------------------------------------------------
%  CUSTOM BOX ENVIRONMENTS
% ----------------------------------------------------------
\usepackage{tcolorbox}
\tcbuselibrary{skins,breakable}

% Info box
\newtcolorbox{infobox}{
  enhanced,
  breakable,
  colback   = softwhite,
  colframe  = navy,
  coltext   = navy,
  boxrule   = 1pt,
  arc       = 4pt,
  left=6pt, right=6pt, top=6pt, bottom=6pt
}

% Warning box
\newtcolorbox{warningbox}{
  enhanced,
  breakable,
  colback   = accentyellow!30,
  colframe  = accentyellow!80!navy,
  coltext   = navy,
  boxrule   = 1pt,
  arc       = 4pt,
  left=6pt, right=6pt, top=6pt, bottom=6pt
}

% Highlight box
\newtcolorbox{highlightbox}{
  enhanced,
  breakable,
  colback   = teamred!10,
  colframe  = teamred,
  coltext   = navy,
  boxrule   = 1pt,
  arc       = 4pt,
  left=6pt, right=6pt, top=6pt, bottom=6pt
}

% ----------------------------------------------------------
%  ITEMIZE / ENUMERATE STYLE
% ----------------------------------------------------------
\usepackage{enumitem}

\setlist[itemize]{
  label=\textcolor{teamred}{\large\textbullet},
  leftmargin=1.2em,
  itemsep=4pt
}

\setlist[enumerate]{
  label=\textcolor{navy}{\arabic*.},
  leftmargin=1.4em,
  itemsep=4pt
}

% ----------------------------------------------------------
%  TITLE STYLE
% ----------------------------------------------------------
\usepackage{titling}

\pretitle{
  \begin{center}
  \headingfont\bfseries\LARGE\color{navy}
}
\posttitle{
  \end{center}
}

\preauthor{
  \begin{center}
  \color{teamred}
}
\postauthor{
  \end{center}
}

% ----------------------------------------------------------
%  OPTIONAL — SOFT PAGE BACKGROUND
%  Comment out if you want white pages
% ----------------------------------------------------------
\usepackage{pagecolor}
\pagecolor{cream}
\color{navy}

% ----------------------------------------------------------
% CUSTOM COMMANDS FOR TEMPLATE
% ----------------------------------------------------------
\newcommand{\courseTitle}{Install Linux On Your Machine}
\newcommand{\courseDate}{25-11-2025}

% ----------------------------------------------------------
% END OF THEME
% ----------------------------------------------------------

\begin{document}

% Title with minimal space
\begin{center}
\headingfont\Huge\bfseries\color{navy}
\courseTitle\\[0.5em]
\Large\color{teamred}
\courseDate
\end{center}

\section*{Topics}

\begin{itemize}
    \item \textbf{Choosing your distro}
          \\Link: \href{https://youtu.be/8yVlJEzq2eg?si=SWcgh_Bw51Li3lv6}{\texttt{https://youtu.be/8yVlJEzq2eg?si=SWcgh\_Bw51Li3lv6}}
    
    \item \textbf{Install your Linux system}
          \\Link: \href{https://www.youtube.com/watch?v=C-a5IamFIuM}{\texttt{https://www.youtube.com/watch?v=C-a5IamFIuM}}
    
    \item \textbf{Switch to Linux while keeping data from Windows / Dual-booting with Windows}
          \\Link: \href{https://youtu.be/XevGfO_vQJQ?si=OBqaM8UlUeH1wjdQ}{\texttt{https://youtu.be/XevGfO\_vQJQ?si=OBqaM8UlUeH1wjdQ}}
    
    \item \textbf{VirtualBox VM setup}
          \\Link: \href{https://youtu.be/wX75Z-4MEoM?si=raCzHOzYJy-GZi9S}{\texttt{https://youtu.be/wX75Z-4MEoM?si=raCzHOzYJy-GZi9S}}
    
    \item \textbf{Post-install (drivers – firewall – backup and restore system)}
          \\Link: \href{https://mmbesar.github.io/Tutorials/Ubuntu-24.04-Post/}{\texttt{https://mmbesar.github.io/Tutorials/Ubuntu-24.04-Post/}}
\end{itemize}

\vspace{1em}
\begin{center}
\color{teamred}\rule{0.8\textwidth}{1pt}
\end{center}

\section*{1. Choosing Your Linux Distribution}

Before installing Linux, you must pick the right \textbf{distro} (distribution). Think of distros like different flavors of the same operating system — they all share the Linux kernel but differ in user interface, package managers, performance, and philosophy.

\subsection*{1.1 What You Should Look For}

When choosing a distro, consider:

\begin{infobox}
\textbf{Distro Selection Guide:}
\begin{itemize}
    \item \textbf{Ease of use} — Beginner-friendly UI, simple updates
    \begin{itemize}
        \item Good Choices: Ubuntu, Linux Mint, Pop!\_OS
    \end{itemize}
    
    \item \textbf{Stability} — Fewer issues \& long-term support
    \begin{itemize}
        \item Good Choices: Ubuntu LTS, Debian
    \end{itemize}
    
    \item \textbf{Performance} — Works well on older hardware
    \begin{itemize}
        \item Good Choices: Xubuntu, Linux Lite
    \end{itemize}
    
    \item \textbf{Security} — Good by default, strong community
    \begin{itemize}
        \item Good Choices: Fedora, Ubuntu
    \end{itemize}
    
    \item \textbf{Software availability} — Access to apps \& drivers
    \begin{itemize}
        \item Good Choices: Ubuntu, Fedora
    \end{itemize}
\end{itemize}
\end{infobox}

\subsection*{1.2 Recommended Distro for Beginners}

\begin{highlightbox}
\textbf{Ubuntu 24.04 LTS}\\
Reason: stable, popular, lots of tutorials, works on almost all laptops.
\end{highlightbox}

\subsection*{1.3 Quick Tips}

\begin{itemize}
    \item If your PC is older than \textbf{10 years}, pick \textbf{Xubuntu} or \textbf{Linux Mint XFCE} or \textbf{MX-Linux} or \textbf{Zorin OS}.
    \item If you're into \textbf{gaming}, choose \textbf{Pop!\_OS} or \textbf{Nobara OS}.
    \item If you want a macOS-like UI, choose \textbf{elementaryOS}.
    \item Avoid Arch-based systems as your first Linux unless you're ready for deeper learning.
\end{itemize}

\section*{2. Installing Your Linux System}

In This section we will use Ubuntu 24.04 for Explanation.

\subsection*{2.1 Before You Install – VERY Important}

\subsubsection*{A. Backup Your Files}

The installer will delete the Windows partition if you're switching fully.
Backup all files from:

\begin{itemize}
    \item Desktop
    \item Documents
    \item Pictures
    \item Music
    \item Videos
\end{itemize}

Move them to:

\begin{itemize}
    \item D:\textbackslash (or another partition)
    \item External HDD
    \item USB stick
\end{itemize}

\subsubsection*{B. Check Your Boot Mode}

Open Start → type \textbf{System Information} → look for:

\begin{itemize}
    \item \textbf{UEFI} → modern, secure boot
    \item \textbf{Legacy BIOS} → older systems
\end{itemize}

This affects partitioning and bootloader behavior.

\subsubsection*{C. Download the ISO}

\begin{itemize}
    \item Ubuntu official website
    \item Size [5 - 6]GB
\end{itemize}

\subsubsection*{D. Create a Bootable USB}

Recommended tool: \textbf{Ventoy}

\begin{itemize}
    \item Supports multiple ISOs on the same flash drive
    \item Very beginner-friendly
\end{itemize}

\subsection*{2.2 Try Ubuntu Before Installing}

Boot from the USB and choose:

\subsubsection*{Try Ubuntu (Live Mode)}

You can test:

\begin{itemize}
    \item Wi-Fi
    \item Bluetooth
    \item Sound
    \item Display
    \item Keyboard
    \item Touchpad
\end{itemize}

\begin{warningbox}
If everything works → safe to install.
\end{warningbox}

\section*{3. Installing Linux (Manual Partitioning)}

This is the \textit{critical} step — mistakes here delete data.

\subsection*{3.1 Choose Manual Installation}

DO NOT choose:

\begin{itemize}
    \item ``Erase disk and install Ubuntu''
    \item ``Install alongside Windows'' (unpredictable for beginners)
\end{itemize}

Instead, choose: \textbf{Something Else (Manual Partitioning)}

\subsection*{3.2 Identify Windows Partitions to Delete}

In the live session, open a terminal:

\begin{highlightbox}
\texttt{\$ lsblk}
\end{highlightbox}

You will see all disks:

\begin{itemize}
    \item Windows system partitions
    \item Data partitions
    \item USB Ventoy
    \item Empty space
\end{itemize}

Delete only:

\begin{itemize}
    \item The Windows partition (C drive)
    \item Windows system reserved partitions
\end{itemize}

\begin{warningbox}
\textbf{Never delete your data partitions.}
\end{warningbox}

\subsection*{3.3 Creating Linux Partitions}

Recommended layout:

\begin{infobox}
\textbf{Linux Partition Layout:}
\begin{itemize}
    \item \textbf{EFI/System} — 512–1024 MB — FAT32 — \texttt{/boot/efi} — Needed only for UEFI
    \item \textbf{ROOT} — 40–100 GB+ — EXT4 — \texttt{/} — Main system
    \item \textbf{SWAP} — Optional (2–4 GB) — swap — Only if RAM < 8GB
    \item \textbf{Data (Windows)} — Keep — NTFS — \texttt{/mnt/data1} — Do \textit{not} format
\end{itemize}
\end{infobox}

\begin{infobox}
\textbf{Beginners Tip:} Don't create a separate \texttt{/home} partition. You can add it later when you're more experienced.
\end{infobox}

\section*{4. Switching to Linux While Keeping Data (Dual Boot)}

\subsection*{4.1 How Dual Boot Works}

Your disk will contain:

\begin{itemize}
    \item Windows NTFS partitions
    \item Linux EXT4 partitions
    \item ESP (bootloader)
    \item Shared data partitions
\end{itemize}

GRUB bootloader will let you choose the OS at startup.

\subsection*{4.2 Before Dual Booting}

\begin{itemize}
    \item Clean Windows (delete temp, uninstall unused programs)
    \item Run: \texttt{chkdsk /f}
    \item Disable Fast Startup in Windows
    \item Shrink Windows partition using \textit{Disk Management}
\end{itemize}

\subsection*{4.3 Partition Strategy}

You will create:

\begin{itemize}
    \item EFI partition (use existing one)
    \item Root partition
    \item (Optional) Swap partition
\end{itemize}

Mount Windows data partitions as:

\begin{highlightbox}
\texttt{/mnt/data1}\\
\texttt{/mnt/data2}
\end{highlightbox}

\begin{warningbox}
\textbf{Never check the ``format'' box.}
\end{warningbox}

\subsection*{4.4 Accessing Windows Files After Installation}

Open \textbf{Files → Other Locations → Ubuntu → mnt}.

Your data partitions appear there.

You can:

\begin{itemize}
    \item Pin them
    \item Bookmark them
    \item Rename them
\end{itemize}

\section*{5. Installing Linux in VirtualBox (VM Setup)}

For students who don't want to touch their main system.

\subsection*{5.1 Why a Virtual Machine?}

\begin{itemize}
    \item Practice Linux commands safely
    \item Test configurations
    \item Use Linux tools on a Windows/macOS system
    \item No risk of data loss
\end{itemize}

\subsection*{5.2 Basic Requirements}

\begin{itemize}
    \item CPU: Dual-Core
    \item RAM: 6–8 GB system → allocate 3–4 GB to VM
    \item Disk: 25–40 GB virtual disk
    \item GPU: Enable 3D Acceleration
\end{itemize}

\subsection*{5.3 Recommended VirtualBox Settings}

\begin{itemize}
    \item \textbf{Chipset:} ICH9
    \item \textbf{EFI:} Enabled for Ubuntu 24.04
    \item \textbf{CPU:} 2 cores
    \item \textbf{Display:} 128 MB VRAM
    \item \textbf{Network:} Bridged Adapter (for networking labs)
\end{itemize}

\subsection*{5.4 Installing Guest Additions}

Inside VM → Devices → Insert Guest Additions CD

Run:

\begin{highlightbox}
\begin{verbatim}
$ sudo apt install build-essential dkms linux-headers-$(uname -r)
$ sudo sh /media/username/VBox*/VBoxLinuxAdditions.run
\end{verbatim}
\end{highlightbox}

This enables:

\begin{itemize}
    \item Auto-resize display
    \item Clipboard sharing
    \item Drag \& drop
    \item Better performance
\end{itemize}

\section*{6. Post-Installation Setup (Ubuntu 24.04)}

\subsection*{6.1 Update System}

\begin{highlightbox}
\texttt{\$ sudo apt update \&\& sudo apt upgrade -y}
\end{highlightbox}

\subsection*{6.2 Install Ubuntu Restricted Extras}

Adds codecs for:

\begin{itemize}
    \item MP3
    \item MP4
    \item MOV
    \item Fonts
\end{itemize}

\begin{highlightbox}
\texttt{\$ sudo apt install ubuntu-restricted-extras}
\end{highlightbox}

\subsection*{6.3 Add Flatpak + Flathub}

\begin{highlightbox}
\texttt{\$ sudo apt install flatpak}\\
\texttt{\$ flatpak remote-add --if-not-exists flathub}\\
\texttt{  https://flathub.org/repo/flathub.flatpakrepo}
\end{highlightbox}

Install Flatpak support in GNOME Software:

\begin{highlightbox}
\texttt{\$ sudo apt install gnome-software gnome-software-plugin-flatpak}
\end{highlightbox}

\subsection*{6.4 AppImage Support}

\begin{highlightbox}
\texttt{\$ sudo apt install libfuse2}
\end{highlightbox}

\subsection*{6.5 GNOME Extensions \& Tweaks}

\begin{highlightbox}
\texttt{\$ sudo apt install chrome-gnome-shell gnome-shell-extension-manager gnome-tweaks}
\end{highlightbox}

\subsection*{6.6 Minimize on Click}

Classic:

\begin{highlightbox}
\texttt{\$ gsettings set org.gnome.shell.extensions.dash-to-dock click-action 'minimize'}
\end{highlightbox}

Newer behavior:

\begin{highlightbox}
\texttt{\$ gsettings set org.gnome.shell.extensions.dash-to-dock click-action 'focus-minimize-or-appspread'}
\end{highlightbox}

\subsection*{6.7 Enable New Document Menu}

\begin{highlightbox}
\texttt{\$ touch \textasciitilde/Templates/new}
\end{highlightbox}

\subsection*{6.8 Keyboard Layout Switching}

Set:
\begin{highlightbox}
\texttt{\$ gsettings set org.gnome.desktop.wm.keybindings switch-input-source "['<Shift>Alt\_L']"}\\
\texttt{\$ gsettings set org.gnome.desktop.wm.keybindings switch-input-source-backward "['<Alt>Shift\_L']"}
\end{highlightbox}

\subsection*{6.9 System Backups}

\subsubsection*{Timeshift (System Snapshots)}

\begin{highlightbox}
\texttt{\$ sudo apt install timeshift}
\end{highlightbox}

\subsubsection*{Déjà Dup (User files backup)}

Built into Ubuntu.

\subsection*{6.10 Final System Update}

\begin{highlightbox}
\texttt{\$ sudo apt update \&\& sudo apt upgrade -y \&\& flatpak update -y \&\& sudo snap refresh}
\end{highlightbox}

\section*{Conclusion}

By the end of this session, students should be able to:

\begin{highlightbox}
\begin{itemize}
    \item \textbf{✔} Choose a Linux distro that fits their hardware and needs
    \item \textbf{✔} Install Linux safely and correctly
    \item \textbf{✔} Dual-boot without losing their Windows files
    \item \textbf{✔} Use Linux inside a Virtual Machine
    \item \textbf{✔} Apply essential post-install steps for performance, security, and usability
\end{itemize}
\end{highlightbox}

% Footer
\vspace{\fill}
\begin{center}
\color{muted}
All rights reserved for \texttt{TogetherTeam} [NOT For Commercial Use] !!\\
contact the author via email \href{mailto:khaledmhfz2004@gmail.com}{\texttt{khaledmhfz2004@gmail.com}}
\end{center}
\end{document}